\paragraph{Theorem F.2.7 (Proved: the orbifold quotient).}
For the specified torus fillings in F.2, with affine
multiplicities $m_1,\ldots,m_k$, the fundamental group of the
total space has a canonical surjection onto
\[
 \langle x_1,\ldots,x_k\mid
       x_i^{m_i}=1,\ x_1\cdots x_k=1\rangle.
 \tag{F.27}
\]
Consequently a simply connected example with $k=2$ requires
$\gcd(m_1,m_2)=1$.  There is no simply connected example with
$k=3$ and $m_1,m_2,m_3\geq2$.  An affine multiplicity-one
point is an ordinary smooth point and reduces to the problem
with one fewer multiple fibre.

\begin{proof}
Over the complement of small discs at the special values,
the torus bundle projects surjectively on fundamental groups
to the punctured base.  Quotient its fundamental group by the
normal closure of all regular-fibre loops.  At a finite
filling its root cover is a smooth torus bundle over a disc,
and its free deck group is cyclic of order $m_i$.
The covering exact sequence shows that the quotient of the
filling group by its fibre subgroup is precisely this cyclic
group.  The boundary base meridian maps to its generator.
Thus adjoining that filling imposes exactly
$x_i^{m_i}=1$ after the fibre loops have been killed.
At the index-one Kuhn cusp the fundamental group is generated
by the three deck loops of its dual torus; these are regular
fibre loops.  Its base meridian contracts.  Hence its filling
imposes the ordinary meridian relation and leaves no extra
generator in this quotient.  The base relation gives the
product relation in (F.27).  Van Kampen proves that (F.27) is
the resulting quotient, in particular the asserted surjection.

For two points, substitute $x_2=x_1^{-1}$; the quotient is
$\mathbb Z/\gcd(m_1,m_2)$.
For three points of orders at least two it is a nontrivial
triangle group.  One can see nontriviality without assuming
a classification of the total space.  Take a geodesic triangle
with angles $\pi/m_1,\pi/m_2,\pi/m_3$ in the sphere, Euclidean
plane or hyperbolic plane according as their sum is greater
than, equal to or less than $\pi$.
In the spherical case the possible orders are, after sorting,
$(2,2,n),(2,3,3),(2,3,4),(2,3,5)$, and the corresponding
spherical triangles exist.  The Euclidean and hyperbolic
angle conditions give the other triangles directly.
If $r_1,r_2,r_3$ are reflections in the opposite sides,
the rotations $r_2r_3,r_3r_1,r_1r_2$ have the stated
orders and product one.  They therefore give a nontrivial
representation of (F.27) into the orientation-preserving
isometry group.  Its quotient cannot be trivial.
Finally a free cyclic action of order one is the identity
construction on a root disc of order one, hence an ordinary
torus bundle filling and not a multiple fibre.
\end{proof}

\paragraph{Proposition F.2.15 (Proved: automatic local complex compatibility).}
Every simply connected datum of F.2 with two finite affine
fillings automatically has finite linear actions commuting
with real complex structures.  This assertion concerns the
individual finite actions and does not produce simultaneous
global periods.

\begin{proof}
Let $m_i$ be the order of the free cyclic affine action
defining the $i$th filling.  Its linear order divides $m_i$.
By F.2.7 the integers $m_1,m_2$ are coprime, so at least one
of them is odd.  The corresponding linear determinant is
one, because a finite real matrix of odd order cannot have
determinant minus one.  The relation
$M_1M_2M_0=I$ and the unipotence of $M_0$ imply
$\det M_1\det M_2=1$, so both determinants are one.

A finite-order real matrix is semisimple.  Its nonreal
eigenvalues occur in conjugate pairs, and its only possible
real eigenvalues are $1$ and $-1$.  Determinant one says
that the $-1$ eigenspace has even real dimension.  Since the
total dimension is six, the $1$ eigenspace also has even
real dimension.  Choose a complex structure on each of
these real eigenspaces by pairing a real basis.  On each
nonreal isotypical real summand, multiplication by $i$ in
its complex eigenspace supplies a commuting complex
structure.  Their direct sum commutes with $M_i$.
The fixed eigenspace is nonzero because the affine action
is free, as in F.1.4, so its even dimension is at least two.

Each chosen structure makes $\mathbb R^6/\mathbb Z^6$ a
complex torus and the given affine map holomorphic; thus it
supplies precisely the local complex-compatibility
hypothesis used in the subsequent reductions.  For three genuine
finite fillings F.2.7 has already excluded simple
connectivity without this argument.
\end{proof}

\paragraph{Proposition F.2.8 (Proved: exclusion of an involutory finite part).}
A simply connected datum with $b_2=0$, $k=2$ and the stated
finite affine fillings has neither finite linear
monodromy of order one or two.  Its unordered pair of finite
linear orders must therefore belong to
\[
 \{(3,4),(3,5),(3,8),(3,10),(4,5),
       (5,6),(5,8),(5,12)\}.
 \tag{F.28}
\]
The affine multiplicities may still be larger than these
linear orders.

\begin{proof}
An identity finite generator would make the product of the
two finite matrices finite, contrary to its nontrivial
unipotent cusp.  Suppose the cohomology generators are
$A,B$ with $A^2=I$, where the other finite linear order is
odd.  It must be odd because the linear orders divide the
coprime affine multiplicities of Theorem F.2.7.
Write $U=AB$; the cusp condition gives $(U-I)^2=0$.
Multiplication of this identity on the left by $A$ and on
the right by $B^{-1}$ gives
\[
 BA+AB^{-1}=2I.
 \tag{F.29}
\]
For $v\in\ker(B-I)$ we obtain $(B+I)Av=2v$.
As $B$ has odd order, $B+I$ is invertible, so $Av=v$.
Every $B$-fixed vector is therefore globally fixed.
The finite free affine action has a nonzero fixed
linear subspace, and F.2.15 makes its real dimension even; it is
at least two.  Consequently the global cohomology invariant
rank is at least two.  The Leray argument of Theorem F.3.1
shows that $b_1=b_2=0$ requires that rank to be exactly one,
a contradiction.  Inverse or transpose monodromy conventions
give the same identity and fixed-space conclusion.

Finally take the complete finite-order list (F.3), discard
one and two, and keep its coprime unordered pairs.
The eight pairs in (F.28) are exactly the result of that
finite enumeration.  This is only a necessary filter, not
a classification of relative positions or affine twists.
\end{proof}

\paragraph{Proposition F.2.9 (Proved: the only first-cohomology pattern for two points).}
For a simply connected two-multiple-fibre datum with $b_2=0$,
the rational
cohomology system $V$ has one invariant, no coinvariants,
and fixed rank two at each finite generator.  Its quotient
$W=V/\mathbb Qe$ by the invariant line has neither invariants
nor coinvariants.  If one finite characteristic polynomial
on $W$ is $(T-1)\Phi_d(T)$ with $\deg\Phi_d=4$, then $W$
is irreducible over $\mathbb Q$.

\begin{proof}
Write $a=\dim V^\Gamma$, $b=\dim(V^*)^\Gamma$ and $f_i$
for the finite fixed ranks.  Theorem F.3.1 gives $a=1$,
vanishing first parabolic cohomology, and
$f_1+f_2=4+b$.  By F.2.15 each $f_i$ is positive and even, so $b$ is
even.  The maximal square-zero cusp has
$\ker N=\operatorname{im}N$.  Therefore every globally
fixed vector belongs to $\operatorname{im}N$, whereas every
globally fixed covector kills $\operatorname{im}N$.
The global invariant and coinvariant spaces pair to zero.

Restriction of a globally invariant covector to the fixed
space of a finite generator is injective: averaging that
generator is the projection onto its fixed space, and the
covector is unchanged by averaging.  Its restriction kills
the $a$ global invariants in that space.  Hence
$a+b\leq f_i$ for each $i$.
If $b\geq2$ is even, the even ranks satisfy
$f_i\geq b+2$, so $4+b=f_1+f_2\geq2b+4$, a contradiction.
Thus $b=0$ and $f_1=f_2=2$.

The quotient $W$ has no invariant covectors, since its dual
embeds into $V^*$.  If it had a fixed line, its preimage in
$V$ would be an extension of the trivial line by the trivial
line.  Every finite generator on this two-dimensional
space would be unipotent of finite order, hence identity.
This would give two invariants in $V$, a contradiction.
Thus $W$ also has no invariant vectors.

If one finite characteristic polynomial is
$(T-1)\Phi_d(T)$ with the second factor irreducible of degree
four, every proper nonzero rational submodule has dimension
one or four.  In dimension one that generator acts as the
identity.  Since the product of the two finite generators
is unipotent, their induced one-dimensional characters are
inverse; hence the other generator also acts trivially.
This would be an invariant line.  In dimension four the
one-dimensional quotient is trivial by the same argument,
giving a coinvariant.  Both have just been excluded.
\end{proof}

\paragraph{Theorem F.2.10 (Proved: exclusion of three further order pairs).}
No simply connected datum with $k=2$, $b_2=0$, and finite linear orders
$(3,5)$, $(3,8)$ or $(3,10)$ has the required maximal
square-zero cusp.  Among such data, the
remaining necessary unordered linear order pairs are
\[
 \{(3,4),(4,5),(5,6),(5,8),(5,12)\}.
 \tag{F.30}
\]
When one order is five, the other finite characteristic
polynomial cannot be $(T-1)^2\Phi_4(T)^2$ or
$(T-1)^2\Phi_6(T)^2$.

\begin{proof}
Let the second order be $d\in\{5,8,10\}$.
Proposition F.2.9 gives an irreducible rational rank-five
quotient $W$, with finite characteristic polynomials
\[
 (T-1)(T^2+T+1)^2,
 \qquad (T-1)\Phi_d(T).
\]
Since the invariant line of $V$ lies in $\operatorname{im}N_V$,
the quotient cusp has Jordan type $(2,2,1)$.
All three local conjugacy classes are self-dual.  Their
centralizer dimensions are respectively $9,5,13$.
For the last value, the conjugate partition is $(3,2)$,
so the dimension is $3^2+2^2$; the finite values are the
sums of squares of their eigenvalue multiplicities.
Consequently the invariant-sheaf Euler formula applied to
$\operatorname{Hom}(W,W^*)$ gives
\[
 \chi\bigl(\mathbb P^1,j_*\operatorname{Hom}(W,W^*)\bigr)
      =-25+9+5+13=2.
 \tag{F.31}
\]
If $W$ were not self-dual, irreducibility would make both
$\operatorname{Hom}_\Gamma(W,W^*)$ and
$\operatorname{Hom}_\Gamma(W^*,W)$ zero.  These are the
degree-zero group and the dual of the degree-two group in
(F.31).  The Euler characteristic would then be nonpositive,
a contradiction.  Thus $W$ has an invariant nondegenerate
rational bilinear form.  Its symmetric part is nonzero,
since an alternating form cannot be nondegenerate in odd
dimension; irreducibility then makes that symmetric part
nondegenerate.  Both generators have determinant one.

Over $\mathbb C$ use
$\operatorname{Spin}_5\simeq\operatorname{Sp}_4$ and its
explicit degree-two representation on the primitive exterior
square.  Choose finite semisimple lifts $R,S$ of the two
orthogonal generators, choosing one of their signs so that
$RS$ is the unipotent lift of their unipotent product.
The orthogonal Jordan type $(2,2,1)$ lifts to a rank-one
transvection.  This can be verified directly, as follows.
For a rank-two square-zero skew-adjoint operator on a
five-dimensional quadratic space, its image is a maximal
isotropic plane, and its nonzero skew form on the dual plane
determines one orbit under the complex orthogonal group.
A rank-one symplectic transvection in dimension four maps,
under the primitive exterior square, to precisely that orbit:
its logarithm has two length-two blocks and one trivial block.
The exact exterior matrices checking these ranks and square
zero are included in \texttt{fourfolds\_order\_pairs.py}.

The first lift has, up to a common sign, eigenvalues
$1,1,\zeta_3,\zeta_3^{-1}$.  Indeed a symplectic torus with
eigenvalues $x,x^{-1},y,y^{-1}$ has primitive exterior
eigenvalues $1,xy,x/y,y/x,(xy)^{-1}$; substitution of the
two equal primitive-third pairs gives the asserted list.
In particular $R$ has a repeated eigenvalue.
The lift $S$ has order dividing $2d$, so it has no eigenvalue
of order three or six.  Nor can it have eigenvalue $1$ or
$-1$: a symplectic pair $x=x^{-1}=\pm1$ would make two of
the nontrivial primitive exterior eigenvalues repeat, whereas
$\Phi_d$ has four distinct roots.
Thus $R^{-1}$ and $S$ have coprime characteristic polynomials.

We finish with an elementary consequence of the rank-one
companion argument.  Two matrices with coprime characteristic
polynomials whose difference has rank one are both cyclic.
To prove this, let $u$ span the difference image.
A common invariant subspace either meets the image nontrivially,
in which case the two induced quotient matrices agree, or is
killed by the difference, in which case their restrictions
agree.  Coprimality forces that subspace to be respectively
the whole space or zero.  The span of the iterates of $u$
under either matrix is stable under both and is nonzero;
it is therefore the whole space, proving cyclicity.
Apply this to $R^{-1},S$, since
$R^{-1}-S=R^{-1}(I-RS)$ has rank one.
But a finite semisimple matrix with a repeated eigenvalue
is not cyclic.  This contradicts the displayed spectrum of
$R$, and excludes all three order pairs.
The same proof gives the final assertion: replace the two
primitive-third pairs on $W$ by two primitive-fourth or
primitive-sixth pairs.  The local centralizer dimension
is still nine, so the Hom Euler characteristic is again two.
The spin lift has, up to sign, eigenvalues
$1,1,\zeta_j,\zeta_j^{-1}$ for $j=4$ or $6$.
The other lift has order dividing ten and has no eigenvalue
$1$ or $-1$ because its primitive exterior spectrum is
simple.  Their spectra are disjoint, and cyclicity gives
the same contradiction.
\end{proof}

\paragraph{Proposition F.2.11 (Proved: rational uniqueness of the extension).}
Let $W$ be a rational rank-five representation of a
twice-punctured affine line, with two finite generators
having fixed rank one and product with unipotent Jordan type
$(2,2,1)$.  Suppose $W$ has neither invariants nor
coinvariants.  Up to rescaling the invariant line, there is
exactly one rational extension
$0\to\mathbb Qe\to V\to W\to0$ in which the two finite
generators retain their finite orders and the cusp has
square-zero logarithm of rank three.  Thus, for the
two-point candidates of F.2.9, the remaining rational
representation problem is exactly the problem of the
rank-five quotients.  This statement does not identify
integral lattices in those representations.

\begin{proof}
Write the two quotient generators as $T_1,T_2$ and the
extension generators in block form
$\left(\begin{smallmatrix}1&c_i\\0&T_i\end{smallmatrix}\right)$.
If $T_i$ has order $m_i$, finite order of the extension is
equivalent to $c_i(I+T_i+\cdots+T_i^{m_i-1})=0$.
Over $\mathbb Q$ this says $c_i\in\operatorname{row}(T_i-I)$,
a space of dimension four.  The two rows give eight
parameters.  Changes of splitting by a row $z\in W^*$
give $c_i\mapsto c_i+z(T_i-I)$; their space has dimension
five because $W^*$ has no invariants.  Hence extensions
split at the two finite cyclic groups form a vector space
of dimension three.

Restriction to the cusp takes this space to
$W^*/\operatorname{row}(T_0-I)$, also of dimension three.
Its kernel is the first parabolic cohomology of $W^*$.
That kernel vanishes: the invariant-sheaf Euler characteristic
is $-5+1+1+3=0$, and its degree-zero and degree-two groups
vanish by the assumed absence of invariants and coinvariants.
Thus cusp restriction is an isomorphism over $\mathbb Q$.

Put $N=T_0-I$ on $W$, so $N^2=0$ and $\operatorname{rank}N=2$.
The cusp extension has nilpotent part
$\left(\begin{smallmatrix}0&c\\0&N\end{smallmatrix}\right)$.
Its square is zero exactly when $cN=0$.
The corresponding subspace of cusp extension classes is
\[
 \frac{\operatorname{ann}(\operatorname{im}N)}
      {\operatorname{row}N},
 \tag{F.32}
\]
which has dimension $3-2=1$.
A class in (F.32) is zero exactly when $c$ vanishes on
$\ker N$; those are precisely the extensions whose cusp
nilpotent part still has rank two.
Every nonzero class raises that rank to three and has square
zero.  Rescaling $e$ acts transitively on the nonzero rational
classes in this one-dimensional space.  This proves both
existence and uniqueness in the stated rational category.
\end{proof}

